% !TeX root = ../main.tex
\ChapterOpening{prologue}{0}{报销主体与凭证要求}
\Subheading{购买方信息与纳税人识别号}
学校报销审核需要将支出与学校业务相联系。开具单位抬头的发票时，\Strong{购买方名称与纳税人识别号应准确填写为中山大学及其对应识别号}。抬头错误、税号遗漏或不匹配，应联系开票方更正，不能仅凭付款记录补足。

铁路电子客票的乘车人实名信息与购买方信息分别承担身份识别和开票用途，\Strong{不能因车票实名而省略单位购买方信息}。航空票据应按实际票种核对填报字段。
\phantomsection\label{invoice-title}
\Subheading{中山大学开票信息}

\begin{PurchaserBox}\Strong{\fontsize{13.5}{20}\selectfont 中山大学 · 发票抬头}\par\vspace{6pt}
\begin{tabularx}{\linewidth}{@{}p{29mm}Y@{}}
\Strong{单位名称} & {\bfseries\fontsize{13.5}{19}\selectfont 中山大学} \\[5pt]
\Strong{纳税人识别号} & {\bfseries\fontsize{13.5}{19}\selectfont 121000004558631445} \\[5pt]
\Strong{地址} & 广东省广州市新港西路135号 \\[5pt]
\Strong{电话号码} & 020-84110285 \\[5pt]
\Strong{开户行} & 中国建设银行广州中山大学支行 \\[5pt]
\Strong{开户行账号} & 44050143004609000001 \\[5pt]
\end{tabularx}
\end{PurchaserBox}
\begin{TeachingPair}
\SmallHeading{购买方与销售方}
购买方信息填写学校名称及识别号；销售方信息由提供商品或服务的商家开具。平台要求填写地址、电话和开户信息时，使用上表对应字段。取得发票后，核对购买方、项目及价税合计。\SideColumn
\begin{SideNote}
\NoteTitle{账号区分}
这里的开户行账号属于学校开票信息。个人收取报销款的银行账户是另一项信息。

\end{SideNote}
\end{TeachingPair}
\Subheading{报销材料的证据链}
完整的报销材料需要支持业务关联、费用内容、实际支付和行程一致性等方面的审核。发票、消费明细、支付记录与行程材料应能相互对应；保单另用于核对被保险人、保障期间和保障范围。

\begin{SummaryBox}
\SummaryTitle{开票合规与业务关联需要分别证明}
抬头正确不代表该项费用当然可以报销。还应结合赛事通知、参赛安排等业务依据，核对费用用途、发生时间及人员。各类凭证的用途与使用场景见第\BookRef{evidence-chain}页。
\end{SummaryBox}

\LessonPage{evidence-chain}{序章}{凭证构成与使用条件}
\Standfirst{按费用性质和审核要求组织凭证。下表说明各类材料的作用，以及需要补充相关材料的常见场景。}
\par\vspace{4pt}{\sffamily\fontsize{9.5}{16}\selectfont
\setlength{\tabcolsep}{5pt}
\renewcommand{\arraystretch}{1.08}
\arrayrulecolor{BookBlue!25}
\begin{tabularx}{\linewidth}{@{}p{27mm}p{24mm}p{44mm}Y@{}}
\toprule
\rowcolor{BookBlue!5}\Strong{\textcolor{BookBlue}{凭证类型}} & \Strong{\textcolor{BookBlue}{凭证性质}} & \Strong{\textcolor{BookBlue}{在报销材料中的作用}} & \Strong{\textcolor{BookBlue}{常见使用场景与注意事项}} \\[5pt]\midrule
\hyperref[invoices]{\Strong{发票／铁路电子客票}}\par\BookRef{invoices} & 交易与开票凭证 & 记载购买方、销售方、项目及金额，为费用入账提供凭证。\Strong{不等同于完税证明，也不能单独证明已付款。} & 电子发票保存平台交付的原文件，PDF、OFD等格式按提交要求使用。截图、裁剪件不能替代电子原文件。 \\[8pt]\midrule
\hyperref[receipts]{\Strong{明细清单（小票）}}\par\BookRef{receipts} & 消费明细 & 解释发票总额的组成，核对商品或服务、数量与单价，识别不属于报销范围的消费。 & 发票仅列笼统项目、需要核对具体消费时补充，例如住宿明细或餐饮清单。\Strong{不能替代正式发票。} \\[8pt]\midrule
\hyperref[payment]{\Strong{支付流水记录}}\par\BookRef{payment} & 支付凭证 & 核对付款人、收款方、金额、时间和支付状态，关联实际垫付与退款情况。 & 需证明实际垫付或解释金额差异时，提供刷卡签购单、银行或支付平台交易详情等。保留成功状态及必要交易信息；是否免附以具体要求为准。 \\[8pt]\midrule
\hyperref[itinerary]{\Strong{行程单／登机牌}}\par\BookRef{itinerary} & 出行明细或乘机证明 & 核对出行人员、日期、路线与参赛安排是否一致。\Strong{不同票种的证明作用不同。} & 网约车行程单用于核对起终点、时间及分段金额。航空行程单和登机牌按实际票种与提交要求留存；预订记录不能单独证明实际乘机。 \\[8pt]\midrule
\hyperref[insurance]{\Strong{保单}}\par\BookRef{insurance} & 保险合同凭证 & 核对被保险人、保障期间、险种、保障范围及保险费。\Strong{保单本身不能证明出行获得学校授权。} & 通常自行提取含关键信息的1～2页；字段跨页时补充必要页面。完整保单留存，保险费发票与付款记录另行整理。 \\[6pt]\bottomrule
\end{tabularx}}\par



\clearpage
\Subheading{发票与铁路电子客票}
\phantomsection\label{invoices}\pdfbookmark[1]{发票}{bm-invoices}
\Standfirst{发票记录交易与开票信息，是费用入账的重要凭证；支付状态、消费明细和业务关联需要相应材料支持。}
\begin{TeachingPair}
\SmallHeading{铁路电子客票}
铁路的\Strong{电子发票（铁路电子客票）}属于全面数字化的电子发票，可用于报销入账。


全国铁路系统已经全面电子化，不再提供纸质车票，所有票据统一从铁路12306平台开具电子发票。在铁路12306 APP下的\strong{“订单”}页面的右上角
点击\strong{“电子发票”}，设置好发票抬头以及接收邮箱后，点击开具电子发票即可选择指定行程开具电子发票。

按铁路12306《业务须知》，乘车日期为\Strong{2025年10月1日及以后}的车票，应在行程结束，或支付退票、改签费用后\Strong{180天内}开具铁路电子客票。
\SideColumn
\begin{SideNote}
\NoteTitle{注意核对开票信息}
抬头与识别号是否正确？金额是否对应？有无退票、改签或重新开票？保留实际有效的凭证。

\end{SideNote}
\end{TeachingPair}
\DocumentFigure[.85\linewidth]{rail-1.png}{0.1}{铁路电子客票示例}
\vspace{2cm}
\SmallHeading{一般电子发票}
除铁路交通之外，其他消费通常\strong{都需要开具电子发票}。电子发票通常由商家或平台提供，常见格式包括 PDF、OFD 等。
请注意，截图、裁剪件或非官方生成的文件不能替代电子原文件。打印提交时尽量使用官方提供的PDF文件直接打印，避免使用截图或裁剪件。



\clearpage
下面是一张代订机票费用的电子普通发票。它显示购买方、销售方、项目与价税合计；购买机票时取得的材料应按实际票种辨认，不能只凭“机票”二字判断。

\DocumentFigure[.94\linewidth]{air-invoice-1.png}{0.2}{电子普通发票}

\begin{SummaryBox}
\SummaryTitle{只需提交发票的PDF}
铁路电子客票提供 PDF 和 OFD 格式，其中 PDF 用于预览和打印。个人提交时只用提交PDF格式的发票即可。
\end{SummaryBox}

\par\vspace{6pt}\begin{minipage}[t]{.31\linewidth}\sffamily\fontsize{9}{16}\selectfont \textcolor{BookPurple}{\Strong{① 购买方}}\par\vspace{3pt}学校名称、识别号与抬头一致。\end{minipage}\hfill\begin{minipage}[t]{.31\linewidth}\sffamily\fontsize{9}{16}\selectfont \textcolor{BookPurple}{\Strong{② 项目}}\par\vspace{3pt}与实际购买的商品或服务对应。\end{minipage}\hfill\begin{minipage}[t]{.31\linewidth}\sffamily\fontsize{9}{16}\selectfont \textcolor{BookPurple}{\Strong{③ 价税合计}}\par\vspace{3pt}需要含税总额，而非税前金额。\end{minipage}\par

\vspace{2cm}

\Subheading{小票}
\phantomsection\label{receipts}\pdfbookmark[1]{小票}{bm-receipts}
小票通常是收银系统、餐饮系统或订单系统生成的消费明细，常列出商品、数量、单价和合计。它适合说明一笔总额由哪些项目组成，
例如一张餐费发票对应的菜品明细。一般只有在采购零售商品时可能会要求提供，但是还是以发票为主。

\begin{TeachingPair}
\Strong{普通收银小票不等同于发票。}只有消费明细，通常还不能替代应取得的正式发票；只有一张笼统的发票，也可能需要小票帮助说明消费内容。

取得小票时，尽量保留商户、日期、明细和合计。热敏纸容易褪色，可及时扫描留存，并按具体费用要求保管纸质原件。\SideColumn
\begin{SideNote}
\NoteTitle{别只看纸张大小}
“小票”是日常叫法。遇到卷式发票等票种，要依据票面名称和法定信息辨认，不能按外形分类。

\end{SideNote}
\end{TeachingPair}

\clearpage
\Subheading{付款记录}
\phantomsection\label{payment}\pdfbookmark[1]{付款记录}{bm-payment}
\par\noindent\begin{minipage}[t]{\dimexpr\linewidth-63mm\relax}\vspace{0pt}
微信、支付宝、银行或购票平台的交易详情，可以帮助说明付款状态、金额、时间及收款方。请选择\Strong{已支付或交易成功}的详情页，
待支付订单不能证明款项已经付出。

付款记录尽量与对应费用关联，如果付款记录和发票对不上可能会被扯皮

\begin{SideNote}
\NoteTitle{发票与付款记录相对应}
图0.2的发票与图0.3的付款记录，金额均为710.00元，商户也可对应。金额相同是核对线索，还应结合交易详情确认属于同一笔费用。

\end{SideNote}\end{minipage}\hfill\begin{minipage}[t]{55.5mm}\vspace{-7pt}
\DocumentFigure{payment-1.png}{0.3}{付款详情}
\end{minipage}\par
\Subheading{行程单}
\phantomsection\label{itinerary}\pdfbookmark[1]{行程单}{bm-itinerary}
网约车行程单通常列出乘车时间、起终点、服务商及各段金额，帮助把交通支出与参赛行程联系起来。
它与发票、付款记录各司其职；若发票合并了多段行程，应核对所含行程和合计金额。

\DocumentFigure{itinerary-1.png}{0.4}{网约车行程单}
\begin{SummaryBox}
\SummaryTitle{注意同名异物}
网约车平台的普通“行程单”通常用于说明行程；航空运输电子客票行程单可能属于正式报销凭证。应结合票面全称、发票标识和开具渠道辨认。

\end{SummaryBox}

\clearpage
\Subheading{保单}
\phantomsection\label{insurance}\pdfbookmark[1]{保单}{bm-insurance}
保单用于说明被保险人、保险期间、保障内容和保险费等信息。完整文件可能连同条款长达\Strong{五六十页}；整理材料时，通常只需自行提取相关的\Strong{1～2页}，让关键信息集中、清楚。

\par\vspace{6pt}\noindent\begin{minipage}[t]{\dimexpr\linewidth-48mm\relax}\vspace{-7pt}
\DocumentFigure{insurance-1.png}{0.5}{保单}
\end{minipage}\hfill\begin{minipage}[t]{42mm}\vspace{0pt}
{\sffamily\fontsize{9}{16}\selectfont \SmallHeading{优先保留的信息：}
\begin{itemize}
\item 保险公司与保单标识
\item 被保险人及身份信息
\item 保险生效、终止日期
\item 险种、保障范围与保险费\end{itemize}
\begin{SideNote}
\NoteTitle{保险费 ≠ 保险金额}
保险费是购买保障所支付的费用；保险金额是约定的保障额度。不要把保额抄成报销金额。

\end{SideNote}
\begin{SideNote}[NotePurple]
\NoteTitle{核对适用范围}
保险期间应当覆盖比赛日期

\end{SideNote}}\end{minipage}\par
\SmallHeading{自行提取 PDF 页面}
打开完整保单，找到包含上述信息的页面，使用 PDF 工具的“提取页面”，或打印时选择相应页码并另存为 PDF。\Strong{保留所选页面的完整内容}，不要只截下姓名或金额；关键字段跨页时，连同必要页面一并提取。

\begin{SummaryBox}
\SummaryTitle{保留完整原件，提交清晰的相关页}
1～2页是常见整理方式，并非固定上限。完整保单留作核对；如具体报销或参赛要求另有规定，按要求补充。

\end{SummaryBox}

\clearpage
\phantomsection\label{material-packaging}
\begin{ReminderBox}
\ReminderTitle{在此衷心建议：}
\begin{ReminderList}
\item 为便于整理，不重不漏地收集票据，强烈建议使用电脑整理报销材料，不要把材料发到文件传输助手再合并转发。
\item 票据采集到之后尽量将他们重命名为可读的文件名(如\\
<time>\_<department>2<destination>\_<vehicle>\_<type>\_<cost>.pdf \\示例： 20261014\_广州2珠海\_高铁车票\_259元.pdf)，
而不是一串数字或者是一串乱码。
\item 统一文件格式，所有票据都转为PDF格式（图片也请转化为PDF格式，请不要发送jpg，png或jpeg等图片格式）。
\item 整理好之后请将所有票据打包成一个压缩包，方便提交。
\end{ReminderList}
\end{ReminderBox}

参考的压缩包格式如下：
\begin{FileTree}
队员甲_<学号>.zip
└── 队员甲_<学号>/
    ├── 穿线/
    │   ├── 穿线发票.pdf
    │   └── 穿线付款记录.pdf
    ├── 打车/
    │   ├── 行程01/
    │   │   ├── 电子发票.pdf
    │   │   └── 电子行程单.pdf
    │   ├── 行程02/
    │   │   ├── 电子发票.pdf
    │   │   └── 电子行程单.pdf
    │   ├── 行程03/
    │   │   ├── 电子发票.pdf
    │   │   └── 电子行程单.pdf
    │   ├── 行程04/
    │   │   ├── 电子发票.pdf
    │   │   └── 电子行程单.pdf
    │   └── 行程05/
    │       ├── 电子行程单.pdf
    │       └── 付款记录.pdf
    ├── 高铁动车/
    │   ├── 铁路电子客票01.pdf
    │   ├── 铁路电子客票02.pdf
    │   └── 铁路电子客票03.pdf
    ├── 机票/
    │   ├── 电子普通发票.pdf
    │   └── 机票付款记录.pdf
    ├── 体检/
    │   ├── 体检发票.pdf
    │   └── 体检付款记录.pdf
    └── 行程报销情况说明.pdf
\end{FileTree}
